\documentclass[a0paper,portrait, blockverticalspace=4em, colspace=2em]{tikzposter}
\tikzposterlatexaffectionproofoff
\usepackage{graphicx}
\usepackage{amsmath}
\usepackage{amssymb}
\usepackage{hyperref}
\usepackage{multicol}
\usetikzlibrary{calc}

% Define ATLAS color scheme
\definecolor{maincolor}{HTML}{9e2b2f}
\definecolor{secondarycolor}{RGB}{158, 43, 47}
\definecolor{lightred}{RGB}{237, 199, 201}

% Special colors for radiation-effect callouts
\definecolor{tid}{RGB}{254,128,2}
\definecolor{niel}{RGB}{0,140,0}
\definecolor{see}{RGB}{190,0,62}
\definecolor{sel}{RGB}{0,90,170}

% ===========================================================================
%  LAYOUT PARAMETERS
% ===========================================================================
\newcommand{\looseitems}{\begin{itemize}\setlength{\itemsep}{0.4em}}
\newcommand{\tightitems}{\begin{itemize}\setlength{\itemsep}{0.2em}}
\newcommand{\captiontext}[1]{{\renewcommand{\baselinestretch}{0.88}\selectfont#1}}

% ---------------------------------------------------------------------------
%  Theme
% ---------------------------------------------------------------------------
\usetheme{Default}
\definecolorstyle{atlascolors}{
    \colorlet{backgroundcolor}{white}
    \colorlet{titlefgcolor}{white}
    \colorlet{titlebgcolor}{maincolor}
    \colorlet{blocktitlefgcolor}{white}
    \colorlet{blocktitlebgcolor}{maincolor}
    \colorlet{blockbodyfgcolor}{black}
    \colorlet{blockbodybgcolor}{lightred!25}
}{}
\usecolorstyle{atlascolors}

% Title bar with ATLAS logo at right
\definetitlestyle{sampletitle}{width=760mm, roundedcorners=20, linewidth=2pt,
  innersep=10pt, titletotopverticalspace=6mm, titletoblockverticalspace=8mm}{%
  \begin{scope}[line width=\titlelinewidth, rounded corners=\titleroundedcorners]
    \draw[color=blocktitlebgcolor, fill=titlebgcolor]
      (\titleposleft,\titleposbottom) rectangle (\titleposright,\titlepostop);
  \end{scope}
  \node[anchor=east, fill=white, rounded corners=10pt, inner sep=10pt, xshift=-22mm]
    at ($(\titleposright,\titlepostop)!0.5!(\titleposright,\titleposbottom)$)
    {
\includegraphics[height=8.6cm]{logos/atlas_transparent.png}};}
\usetitlestyle{sampletitle}

\title{\parbox{0.74\linewidth}{\centering\huge
    Upgrade of the ATLAS Tile Calorimeter Front-End\\[2mm]Power Supply for the High-Luminosity LHC}}
\author{Seyedali Moayedi \quad on behalf of the ATLAS Tile Calorimeter System}
\institute{University of Texas at Arlington}
\date{}

\setlength{\columnsep}{0.5cm}

\begin{document}
\maketitle

\begin{columns}

% ===========================================================================
%  COLUMN 1 
% ===========================================================================
\column{0.333}

\block{The ATLAS Tile Calorimeter}{
    \looseitems
        \item As the primary hadronic sampling calorimeter of ATLAS, TileCal is crucial for measuring hadronic jets, missing transverse energy, hadronically decaying $\tau$-leptons, and for muon identification.
        \item In total, there are 256 modules in the ATLAS Tile Calorimeter.
        \item To cope with the unprecedented instantaneous luminosity and pile-up of the High-Luminosity LHC (HL-LHC), the entire on-detector electronics system is being fully replaced.
    \end{itemize}
}

\block{LVPS \& HL-LHC Upgrade}{
    \looseitems
        \item The Low Voltage Power Supplies (LVPS) are located inside a box in the ``finger'' of each module and provide power to front-end electronics to read out PMT signals.
        \item The LVPS bricks are switching DC-DC converters with a topology of two-switch forward converter and a switching frequency of 300\,kHz.
        \item In the Phase-II upgrade system, all bricks will convert a 200\,V input to a nominal 10\,V output at a nominal power of 23\,W. Point of load regulators on front-end electronics provide different necessary voltages.
    \end{itemize}
}

\block{LVPS Brick Segmentation}{
    \looseitems
        \item In the legacy system, 8 different bricks inside the box provided power to each module at 8 different voltage levels.
        \item An improvement of the new LVPS design is the segmentation of the power supply into 8 independent, identical bricks providing a single voltage.
        \item This topology enhances the reliability of the readout systems and minimizes the impact of a possible failure during data taking.
    \end{itemize}
    \vspace{0.3cm}
    \begin{center}
        \includegraphics[width=1.0\linewidth]{assets/fig_lvps_brick.pdf}
    \end{center}
    \vspace{-0.2cm}
    \captiontext{The new Tile Calorimeter LVPS brick design, transforming 200\,V DC to a steady 10\,V DC output for local regulation.}
}

\block{V74X Parameter Specifications}{
    The target performance parameters for the V74X LVPS bricks are strictly controlled to ensure sufficient thermal dissipation and voltage stability.
    \vspace{0.4cm}
    
    \begin{center}
    \renewcommand{\arraystretch}{1.3}
    \begin{tabular}{|l|l|}
    \hline
    \textbf{Parameter} & \textbf{Specification} \\
    \hline
    Input Voltage & 200\,V DC \\
    Output Voltage & 10\,V (nominal) \\
    Nominal Power & 23\,W \\
    Topology & Two-switch forward \\
    Switching Frequency & 300\,kHz \\
    Target Efficiency & $\geq$ 72\% \\
    Cooling Plant Cap. & 77\,kW \\
    \hline
    \end{tabular}
    \end{center}
}

\block{References}{
    \small
    \begin{enumerate}\setlength{\itemsep}{0.2em}
        \item ATLAS Collaboration, \emph{Technical Design Report for the Phase-II Upgrade of the ATLAS Tile Calorimeter}, CERN-LHCC-2017-019, ATLAS-TDR-028 (2018).
        \item S. Moayedi (ATLAS), \emph{Upgrade of the ATLAS Tile Calorimeter Front-End Power Supply for the HL-LHC}, TWEPP 2025.
    \end{enumerate}
}

% ===========================================================================
%  COLUMN 2 
% ===========================================================================
\column{0.333}

\block{Radiation Criteria \& Test Facilities}{
    \tightitems
        \item In the Tile Calorimeter, the fingers are the harshest location in terms of radiation. Early brick designs showed excessive trips, mandating a customized rad-hard design.
        \item Based on new guidelines, batch-to-batch safety factors are not applicable; each batch used on final boards must be tested for \textcolor{tid}{TID}, \textcolor{niel}{NIEL}, and \textcolor{see}{SEE}.
    \end{itemize}
    \vspace{0.3cm}
    \begin{center}
        \includegraphics[width=1.0\linewidth]{assets/fig_lvps_tests_overview.pdf}
    \end{center}
    \vspace{-0.2cm}
    \captiontext{Overview of irradiation test facilities (CHARM, JSI, PSI, PIF, CC60) and respective particle types and energies.}
}

\block{Tested Analog \& Logic Components}{
    \tightitems
        \item \textbf{Diodes:} Seven diodes (MURS160T3G, VS-8ETH06, SMCJ15A, etc.) were tested. The TVS diode (SMCJ15A) \textcolor{tid}{TID} tests were successfully performed by biasing at stand-off voltage during beam time at 2.85\,Gy/h.
        \item \textbf{Op-Amp \& Voltage Reg:} LTC6241 and LT3080 showed no significant changes in parameters during \textcolor{see}{SEE} and \textcolor{tid}{TID} tests.
        \item \textbf{MOSFET Driver:} IR2110 was monitored for \textcolor{see}{SEEs} up to $3.4 \times 10^{11}$\,p/cm$^2$ and showed no events. \textcolor{tid}{TID} tests at 3\,Gy/h showed no significant parameter changes.
    \end{itemize}
}

\block{Design Improvements \& Efficiency}{
    \tightitems
        \item Efficiency and thermal management were limiting factors in earlier revisions, requiring at least 65\% efficiency for the 77\,kW cooling plant capacity.
        \item The original power MOSFETs (STB57N65M5) caused high switching losses. They were replaced by SIHFS9N60A, which features a 65\% smaller gate-source capacitance.
        \item This replacement resulted in a significant efficiency increase from $\sim$58\% to $\sim$72\%, easing cooling constraints.
        \item Additionally, the output inductor was replaced with a smaller series resistance version to reduce localized heating.
    \end{itemize}
}

\block{LVPS Tested Components Matrix}{
    Radiation performance matrix for the major LVPS components during qualification. Components must pass all metrics to be selected for the final design.
    \vspace{0.4cm}

    \begin{center}
    \renewcommand{\arraystretch}{1.3}
    \begin{tabular}{|l|l|c|c|c|l|}
    \hline
    \textbf{Type} & \textbf{Component} & \textbf{\textcolor{tid}{TID}} & \textbf{\textcolor{niel}{NIEL}} & \textbf{\textcolor{see}{SEE}} & \textbf{Status} \\
    \hline
    Diode & MURS160T3G & Pass & Pass & Pass & Selected \\
    Diode & VS-8ETH06 & Pass & Pass & Pass & Selected \\
    MOSFET & SIHFS9N60A & Pass & Pass & Pass & Selected \\
    MOSFET Drv & IR2110 & Pass & Pass & Pass & Selected \\
    Op-Amp & LTC6241 & Pass & Pass & Pass & Selected \\
    Volt Reg & LT3080 & Pass & Pass & Pass & Selected \\
    Isol. Amp & SI8920 & Pass* & Pass & Pass* & Selected \\
    Isol. Amp & HCPL7800 & Pass & Fail & Pass & Rejected \\
    \hline
    \end{tabular}
    \end{center}
    \vspace{0.2cm}
    {\small *Pass contingent on maintaining an input voltage $\leq 100$\,mV.}
}

% ===========================================================================
%  COLUMN 3 
% ===========================================================================
\column{0.333}

\block{Wire Reduction \& Control Signals}{
    \looseitems
        \item The number of signals required to control each brick had to be aggressively reduced due to limitations in the long cabling running from the service room to the detector modules.
        \item The logic for controlling the brick on/off state was consolidated to improve reliability.
    \end{itemize}
    \vspace{0.3cm}
    \begin{center}
        \includegraphics[width=1.0\linewidth]{assets/fig_wire_reduction.pdf}
    \end{center}
    \vspace{-0.2cm}
    \captiontext{Schematic of the simplified control logic minimizing cable plant overhead.}
}

\block{Isolation Amplifier Radiation Tests}{
    \looseitems
        \item \textbf{SI8920:} Initially showed fluctuations in output voltage during \textcolor{see}{SEE}/\textcolor{tid}{TID} tests at $3.5 \times 10^8$\,p/cm$^2$/s when operated with 150\,mV input. Further testing at CHARM and PSI demonstrated stable behavior restricted to input voltages below 100\,mV.
        \item \textbf{HCPL7800:} Tested as an alternative at $2.1 \times 10^8$\,p/cm$^2$/s. Small Single Event Transients (\textcolor{see}{SETs}) of $\sim$100\,mV for 5\,$\mu$s were filtered by the brick's low-pass circuits, passing \textcolor{see}{SEE} and \textcolor{tid}{TID}. However, \textcolor{niel}{NIEL} tests revealed failure due to displacement damage at $8 \times 10^{12}$\,n/cm$^2$.
        \item \textbf{Resolution:} The SI8920 was definitively chosen for the production design, with brick circuits modified to cap its input voltage at 100\,mV.
    \end{itemize}
    \vspace{0.12cm}
    \begin{center}
        \includegraphics[width=1.0\linewidth]{assets/fig_si8920_isolation_amp.pdf}
    \end{center}
    \vspace{-0.2cm}
    \captiontext{Test results for the SI8920 isolation amplifier under proton irradiation.}
}

\block{System Level Tests (CHARM)}{
    \looseitems
        \item Pre-production bricks (utilizing the HCPL7800) were tested in the mixed field of the CHARM facility to validate system-level resilience before mass production.
        \item The legacy power MOSFETs failed at a dose of 323\,Gy (far exceeding the 80\,Gy criterion) by shifting the gate-source threshold to zero, causing a short-circuit.
        \item The HCPL7800 failed at $2.75 \times 10^{12}$\,n/cm$^2$ in the mixed field, reinforcing its replacement.
    \end{itemize}
}

\block{Conclusions \& Outlook}{
    \looseitems
        \item The LVPS brick design has been heavily optimized for the HL-LHC radiation environment, ensuring operational stability and efficiency $\geq$ 72\%.
        \item The final design leveraging the SI8920 and SIHFS9N60A meets all HL-LHC requirements and is ready for production.
        \item Around 930 DBs and associated power boards will be produced during 2026 and 2027 to fulfill the upgrade timeline.
    \end{itemize}
}

\end{columns}
\end{document}
